\section{支持脚トルク低減効果と制御的解釈}
\label{sec:discussion_torque}

支持脚ロール方向トルクの挙動を見ると，Case4 における提案手法適用時には，DS から SS へ切り替わる瞬間に生じるトルクのピーク値が明確に低減していた．これは，CoM 高さの増大によって横方向の倒れ込みが緩和され，支持脚にかかるトルクが低減した結果であると解釈できる．

この結果は，提案手法が単に ZMP の位置を制御するだけでなく，支持脚関節に作用する力学的負荷を低減する効果を持つことを示している．したがって，本手法は歩行安定性の向上に加え，アクチュエータ負荷の軽減という観点からも有用であると考えられる．